\chapter{Kajian Pustaka}
\label{bab:kajian-pustaka}

Awali pembahasan pada bab ini dengan mengutarakan \emph{Problem} pada
\emph{Problem Domain} dari tesis, diiringi teori-teori penunjang yang terkait.
Untuk menguatkan klaim orisinalitas, berikan ulasan tentang penelitian terkait
yang juga mencoba menyelesaikan \emph{Problem} tersebut.

\section{Teori Penunjang}

Uraikan dasar teori, temuan, dan bahan dari pustaka ilmiah lain yang dijadikan
landasan untuk melakukan tesis yang diusulkan \cite{pengarang2014buku}.

\section{Penelitian Terkait}

Penelitian terkait memuat hasil penelitian pihak lain yang mempunyai
\emph{Problem} yang sama, tetapi dengan \emph{Uniqueness} yang berbeda. Penulis
harus mencantumkan sumber referensi secara jujur dan benar seperti ini
\cite{pengarang2013jurnal}. Pencantuman sumber referensi dilakukan baik terhadap
kutipan langsung maupun kutipan tidak langsung (parafrase).

Untuk kutipan langsung dan pendek (1--2 baris), gunakan \verb|\kutipan{...}|,
seperti ini \kutipan{Ini contoh penulisan rujukan untuk kutipan langsung dan
pendek} \cite{pengarang2013jurnal,pengarang2012seminar}. Untuk kutipan langsung
dan panjang (lebih dari 2 baris), gunakan lingkungan \texttt{kutipanpanjang}
dengan rujukan sebagai argumen opsional:

\begin{kutipanpanjang}[\cite{mahasiswa2011tesis}]
Ini contoh penulisan rujukan untuk kutipan langsung dan panjang, ditulis pada
paragraf terpisah dengan memberikan tanda petik ganda di awal dan akhir kutipan,
ukuran font 10 point dan margin kanan kiri yang masuk 10 mm dari batas penulisan,
kemudian diiringi dengan sumber referensi pada daftar pustaka.
\end{kutipanpanjang}

Untuk kutipan tidak langsung (parafrase), penulis dapat menuliskan sumber
referensi setelah kalimat parafrase selesai, seperti ini \cite{penulis2012media}.
